\section{结论}

本文提出了一种面向结构化分析任务的数据感知多智能体规划框架，并在专利分析场景中构建了独立于学位论文的期刊稿证据工作区。与只优化执行层的系统不同，本文把关注点前移到规划层，并通过 Evidence Grounding、Knowledge-Constrained Planning 与 Iterative Refinement 三项机制，把数据证据、知识约束和执行反馈组织到统一闭环中。

基于已冻结的 90 组内部重复实验与已完成的 30 组外部机制基线，本文已经可以得到三项稳健判断：其一，数据感知是推动规划从问题语义驱动走向证据驱动的关键机制；其二，两轮迭代决定了系统可达到的性能上限；其三，知识图谱在平均意义上提供正向贡献，但收益强度受问题结构显著影响。换言之，完整系统当前在内部机制归因层面已经表现出“总体最优且相对稳定”的特征，而 \texttt{B\_data\_aware} 则稳定优于不含数据感知的方案。

\IfFileExists{generated/table_external_llm.tex}{
外部机制基线已经说明，仅靠单智能体 ReAct 风格规划或仅靠执行反馈，并不足以达到完整系统的稳定表现。接下来的增强版任务是把这一判断扩展到五问题口径、子领域稳定性、Bugzilla 跨领域小验证和人工盲评，从而形成“内部归因 + 外部机制对比 + 外部有效性”的完整闭环。若这些补充证据与当前趋势保持一致，则本文可以进一步支持这样一个结论：本文提出的是一种具有外部可接受性的规划方法，而不只是针对单一评分器优化的实验系统。
}{
外部机制基线、子领域稳定性、跨领域验证和人工盲评接口已经接入期刊稿工作区，但若当前编译版本尚未显示对应表图，则相关证据仍在回填过程中。对此，本文保持审慎口径：在外部证据真正落地之前，不把“完整系统在外部对比和外部有效性验证中均最优”写成既成事实。
}

总体而言，本文把结构化分析系统的核心问题从“如何把代码跑出来”前移到“如何先形成更合理的分析计划”，并给出了一条可复验、可扩展、可逐层验证的软件方法路线。后续工作将继续围绕五问题扩展、相邻子领域稳定性重跑、Bugzilla 跨领域验证和更大规模人工盲评完善证据链，同时探索知识约束强度的自适应调节机制。
